% AY2024 制御工学 解答例
% 回答の流れ・言葉遣いは 03_制御工学/参考/「制御工学Ⅱ（完成版）」に寄せる。
% デザイン（囲み枠等）は真似しない。解答・数値は本年度過去問から自力で導いた。
\documentclass[11pt,a4paper]{article}
\usepackage{../../../00_common/preamble}

\usetikzlibrary{positioning,arrows.meta,calc}
\usepackage{pgfplots}
\pgfplotsset{compat=newest}

\title{制御工学 AY2024 解答例（専門応用科目I）}
\author{}
\date{}

\begin{document}
\maketitle

% ==================================================================
\section*{第1問〔I〕}

壁 $-$（バネ $k_2$ ∥ ダンパ $d$）$-$ 質量 $m$ $-$ バネ $k_1$ $-$ 入力端 A（変位 $x_1$）の系。

\subsection*{(1) 運動方程式と伝達関数}
$k_2,d$ は質量と壁を、$k_1$ は質量と入力端 $x_1$ を結ぶので
\begin{align*}
  &\qquad m\ddot{x}_2=-k_2 x_2-d\dot{x}_2+k_1(x_1-x_2)\\
  &\therefore\quad m\ddot{x}_2+d\dot{x}_2+(k_1+k_2)x_2=k_1 x_1
   \quad\therefore\quad (ms^2+ds+k_1+k_2)X_2=k_1 X_1.
\end{align*}
\[
  \therefore\quad \ans{\;G(s)=\dfrac{X_2(s)}{X_1(s)}=\dfrac{k_1}{ms^2+ds+(k_1+k_2)}\;}
\]

\subsection*{(2) ステップ応答（$m=1,d=2,k_1=3,k_2=2$）}
$G(s)=\dfrac{3}{s^2+2s+5}$、極は $s=-1\pm2j$（不足減衰）。$X_1=\dfrac1s$ より
\begin{align*}
  &\qquad X_2(s)=\frac{3}{s(s^2+2s+5)}=\frac{3/5}{s}-\frac{3}{5}\cdot\frac{s+2}{(s+1)^2+4}\\
  &\therefore\quad x_2(t)=\frac35-\frac35\e^{-t}\!\left(\cos2t+\tfrac12\sin2t\right).
\end{align*}
$G(0)=\dfrac35$ より定常値は $\dfrac35$。
\[
  \therefore\quad \ans{\;x_2(t)=\frac35-\frac35\e^{-t}\!\left(\cos2t+\tfrac12\sin2t\right),\qquad x_2(\infty)=\frac35\;}
\]

\subsection*{(3) 最大値と発生時間}
\begin{align*}
  &\qquad \dot{x}_2(t)=\frac32\e^{-t}\sin2t=0
   \quad\therefore\quad \sin2t=0.
\end{align*}
最初のピークは $2t=\pi$、すなわち $t=\dfrac{\pi}{2}$。このとき $\e^{-t}(\cos2t+\tfrac12\sin2t)=\e^{-\pi/2}(\cos\pi)=-\e^{-\pi/2}$ なので
\[
  \therefore\quad \ans{\;x_{2\max}=\frac35\bigl(1+\e^{-\pi/2}\bigr)\approx0.72\ \ \Bigl(t=\tfrac{\pi}{2}\approx1.57\Bigr)\;}
\]

\subsection*{(4) 無減衰の場合（$m=1,d=0,k_1=3,k_2=1$）}
\begin{align*}
  &\qquad G(s)=\frac{3}{s^2+4}
   \quad\therefore\quad X_2(s)=\frac{3}{s(s^2+4)}=\frac34\!\left(\frac1s-\frac{s}{s^2+4}\right)\\
  &\therefore\quad x_2(t)=\frac34\bigl(1-\cos2t\bigr).
\end{align*}
\[
  \therefore\quad \ans{\;x_2(t)=\dfrac34(1-\cos2t)\;}
\]
$d=0$ なので減衰せず、$\dfrac34$ を中心に振幅 $\dfrac34$ で持続振動する。

\subsection*{(5) A を固定端、質量にステップ外力（$m=1,d=2,k_1=1,k_2=1$）}
A を固定すると両バネとも壁につながり、合成ばね定数は $k_1+k_2=2$。運動方程式は
\[
  m\ddot{x}_2+d\dot{x}_2+(k_1+k_2)x_2=f(t)\ \Longrightarrow\ \ddot{x}_2+2\dot{x}_2+2x_2=f(t).
\]
$f=\dfrac1s$（単位ステップ）、極 $s=-1\pm j$ より
\begin{align*}
  &\qquad X_2(s)=\frac{1}{s(s^2+2s+2)}=\frac{1/2}{s}-\frac12\cdot\frac{s+2}{(s+1)^2+1}\\
  &\therefore\quad x_2(t)=\frac12-\frac12\e^{-t}(\cos t+\sin t).
\end{align*}
\[
  \therefore\quad \ans{\;x_2(t)=\dfrac12-\dfrac12\e^{-t}(\cos t+\sin t)\;}\qquad(\text{定常値}\ \tfrac12)
\]

\bigskip
\noindent 検算: (2) 定常 $\tfrac35$・極 $-1\pm2j$、(3) $t=\tfrac\pi2$ で $x_{2\max}=\tfrac35(1+\e^{-\pi/2})$、
(4) 無減衰振動、(5) 定常 $\tfrac12$ を SymPy で確認✓。

\newpage
% ==================================================================
\section*{第2問〔II〕}

\subsection*{(1) 図2 一巡伝達関数のゲイン線図（$G_c=\frac1s,\ G_p=\frac{s+1}{10s+1}$）}
一巡伝達関数は
\[
  L(s)=G_c G_p=\frac{1}{s}\cdot\frac{s+1}{10s+1}=\frac{s+1}{s(10s+1)}.
\]
積分器 $\frac1s$ により低周波は $-20\,\mathrm{dB/dec}$（$|L|\approx\frac1\omega$、$\omega=1$ で $0\,\mathrm{dB}$ を通る漸近線）。
折れ点は $10s+1$ より $\omega=0.1$（極、$-20$ 追加）、$s+1$ より $\omega=1$（零点、$+20$ 追加）。
傾きは $\omega<0.1$ で $-20$、$0.1<\omega<1$ で $-40$、$\omega>1$ で $-20\,[\mathrm{dB/dec}]$。

\begin{center}
\begin{tikzpicture}
\begin{axis}[
  width=0.86\linewidth, height=5.2cm,
  xmin=1e-2, xmax=1e2, ymin=-70, ymax=50, xmode=log,
  xlabel={$\omega\,[\mathrm{rad/s}]$}, ylabel={ゲイン [dB]},
  ytick={-60,-40,-20,0,20,40},
  xmajorgrids=true, ymajorgrids=true, xminorgrids=true,
  grid style={line width=0.3pt, draw=gray!60},
  extra y ticks={0}, extra y tick style={grid=major, grid style={line width=1pt, draw=black}},
  legend pos=north east, clip=true, enlargelimits=false,
]
\addplot[thick] coordinates {(0.01,40) (0.1,20) (1,-20) (100,-60)};
\addlegendentry{漸近線}
\addplot[thick,dashed] coordinates
  {(0.01,40) (0.1,17) (0.3,1) (1,-17) (10,-40) (100,-60)};
\addlegendentry{真値}
\node[anchor=south west] at (axis cs:0.1,20) {\scriptsize $\omega=0.1$};
\node[anchor=north east] at (axis cs:1,-20) {\scriptsize $\omega=1$};
\end{axis}
\end{tikzpicture}
\end{center}

\subsection*{(2) 図2 のステップ応答}
閉ループは
\begin{align*}
  &\qquad T(s)=\frac{G_cG_p}{1+G_cG_p}
   =\frac{s+1}{s(10s+1)+(s+1)}=\frac{s+1}{10s^2+2s+1}.
\end{align*}
開ループが積分器をもつ（1 型）ので、ステップ定常偏差は $0$、$T(0)=1$。極は $s=-0.1\pm0.3j$。
$X(s)=\dfrac{T(s)}{s}$ を逆変換すると
\[
  \therefore\quad \ans{\;y(t)=1-\e^{-0.1t}\cos0.3t\;}
\]
（緩やかに減衰する振動で、定常値 $1$ に収束する。）

\subsection*{(3) 図3 の伝達関数（$G_c=\frac{1}{K_1},\ G_p=\frac{1}{K_2}$）}
各加え合わせ点・積分器の関係（$G_c$ 出力 $\to$ sum1 へ、2 段目積分器出力 $\to$ sum2 へ、
出力 $Y\to$ sum3 へ、いずれも負帰還）を連立して整理すると
\[
  \therefore\quad \ans{\;\dfrac{Y(s)}{U(s)}=\dfrac{1}{K_1K_2\,s^2+(K_1+2K_2)\,s+1}\;}
\]

\subsection*{(4) 図3 のステップ応答（$K_1=\frac12,\ K_2=\frac13$）}
$K_1K_2=\dfrac16$、$K_1+2K_2=\dfrac12+\dfrac23=\dfrac76$ より
\begin{align*}
  &\qquad T(s)=\frac{1}{\frac16 s^2+\frac76 s+1}=\frac{6}{s^2+7s+6}=\frac{6}{(s+1)(s+6)}.
\end{align*}
極は実数 $s=-1,-6$（過減衰）、$T(0)=1$。$X(s)=\dfrac{T(s)}{s}$ を分解して
\[
  \therefore\quad \ans{\;y(t)=1-\dfrac65\e^{-t}+\dfrac15\e^{-6t}\;}\qquad(\text{定常値}\ 1)
\]

\bigskip
\noindent 検算: (2) $T=(s+1)/(10s^2+2s+1)$・$y=1-\e^{-0.1t}\cos0.3t$・定常 $1$、
(3) 連立解 $T=1/(K_1K_2 s^2+(K_1+2K_2)s+1)$、(4) $T=6/((s+1)(s+6))$ を SymPy で確認✓。

\end{document}
